% Original explanation prompted by Rajshahi 2025 CQ; technical basis: MDN server-side introduction.
\subsubsection{ফ্রন্ট-এন্ড ও ব্যাক-এন্ড: দৃশ্যমান অংশ ও তথ্যের কাজ}
ওয়েবসাইট তৈরির কাজকে সাধারণভাবে \textbf{ফ্রন্ট-এন্ড (Front-end)} ও \textbf{ব্যাক-এন্ড (Back-end)} হিসেবে বোঝানো হয়। ফ্রন্ট-এন্ড হলো ব্যবহারকারী ব্রাউজারে যে পেজ, মেনু, বোতাম ও ফর্ম দেখে এবং ব্যবহার করে তার অংশ। HTML তথ্যের গঠন, CSS বিন্যাস এবং JavaScript ব্যবহারকারীর সঙ্গে মিথস্ক্রিয়ায় সাহায্য করে। ব্যাক-এন্ডে সার্ভারের প্রোগ্রাম অনুরোধ প্রক্রিয়া করে, তথ্য যাচাই করে এবং প্রয়োজন হলে ডেটাবেজ থেকে তথ্য আনে বা সংরক্ষণ করে।

\begin{keypointbox}{কলেজের ফলাফল দেখার উদাহরণ}
রোল লেখার ঘর ও ফল দেখানোর পেজ ফ্রন্ট-এন্ড। দেওয়া রোলের ফল খুঁজে পাঠানোর সার্ভারের কাজ ব্যাক-এন্ড। শুধু HTML দিয়ে ঘর তৈরি করলেই ফল খোঁজার ব্যবস্থা তৈরি হয় না। একটি সাধারণ স্থির নোটিশ পেজে আলাদা ডেটাবেজ বা নিজস্ব ব্যাক-এন্ড প্রোগ্রাম লাগতেই হবে এমন নয়।
\end{keypointbox}
ফ্রন্ট-এন্ড ও ব্যাক-এন্ড ওয়েবসাইট উন্নয়নের দুটি কাজের দিক। ক্লায়েন্ট ও সার্ভার অনুরোধ ও সাড়া আদান-প্রদানের দুটি ভূমিকা। এই দুই পার্থক্যের ভিত্তি আলাদা।

